\section{Results}\label{sec:results}
We first compare the filter-trained policies with the alternatives under exact battery parameters. Table~\ref{tab:viol} gives the policies trained from scratch and Table~\ref{tab:ws} the warm-started ones; Tables~\ref{S-tab:perseed} and~\ref{S-tab:perseedws} of the supplementary material give them per seed, and Section~\ref{sec:ablation} varies one factor at a time.

\subsection{Main comparison}\label{sec:main}
Without the filter, the baseline policies violate both bus limits at every state of charge (Table~\ref{tab:viol}). Trained from scratch, at state of charge 0.2 they spend \res{base_uv_pct_soc0.2}\% of the episode below $V_{\min}$ and \res{base_oc_pct_soc0.2}\% above $I_{\lim}$. The warm-started baseline policies, trained for \res{train_steps_ws} further steps, behave alike (Table~\ref{tab:ws}).

The warm-started filter-trained policies never leave either limit, never collapse the bus and never fall. At every state of charge they hop more often than the warm-started baseline policies, but lower (Table~\ref{tab:ws}).
Figures~\ref{fig:traces}, \ref{fig:violsoc} and~\ref{fig:tradeoff} show one episode without and with the filter, the violation times against state of charge, and the max height of the warm-started policies, grouped into those that violate a limit and those within both limits. At every state of charge, each method without violation time has a lower mean max height than every baseline and penalty policy that violates a limit.

The penalty policies trained from scratch ($w=\res{pen_w_scratch}$) have no violation and no hops (Table~\ref{tab:viol}).
\begin{table}[tbp]\centering\caption{Policies trained from scratch: violations and task performance by method and state of charge (SOC) (mean over seeds; \res{eval_eps_main} episodes per seed and state of charge). Falls: fraction of episodes. Standard deviations, lowest bus voltage, mean apex and return: full version, Table~\ref{S-tab:violfull} of the supplementary material.}\label{tab:viol}
\gentable{tab_violations_short}
\end{table}
\begin{table}[tbp]\centering\caption{Warm-started policies with three seeds, each continued from the baseline policy of the same seed: violations and task performance by method and state of charge (mean over seeds; \res{eval_eps_main} episodes per seed and state of charge). Falls: fraction of episodes. The full version, with standard deviations, the penalty weights \res{pen_ws_w_part} (one seed) and the other policies with the filter at test time, is Table~\ref{S-tab:wsfull} of the supplementary material.}\label{tab:ws}
\gentable{tab_warmstart_short}
\end{table}
\begin{figure}[tbp]\centering\genfig{F2_traces}{\linewidth}\caption{One episode at state of charge 0.2 (reset seed 0) of the warm-started policies in baseline mode (baseline\_ws\_s0) and filter mode (filter\_ws\_s0). From top: rail height above its start; lowest bus voltage and highest battery current within each \res{env_dtctrl_ms}~ms control step, with the limits $V_{\min}$ and $I_{\lim}$; fraction of substeps in which the filter is active.}\label{fig:traces}\end{figure}
\begin{figure}[tbp]\centering\genfig{F3_violation_vs_soc}{\linewidth}\caption{Time per episode below $V_{\min}$ and above $I_{\lim}$ against state of charge, for the warm-started policies in their own training modes. Markers: mean over seeds; bars: minimum to maximum over seeds; hollow markers and dashed lines: methods with one seed; black rings labelled ``fell $k/n$'' with the method: $k$ of the $n$ episodes (over seeds) ended in a fall (violation times do not show falls).}\label{fig:violsoc}\end{figure}\begin{figure}[tbp]\centering\genfig{F6_tradeoff}{\linewidth}\caption{Max height of the warm-started policies in the main evaluation, one panel per state of charge: rows that violate a limit (time below $V_{\min}$ or above $I_{\lim}$; hollow markers) and rows within both limits (no such time in any episode; filled markers). Markers: mean over \res{n_seeds} seeds; bars: standard deviation over seeds; \res{eval_eps_main} episodes per seed and state of charge; ``fell $k/n$'': $k$ of the $n$ episodes over seeds ended in a fall. The penalty weights \res{pen_ws_w_part} (one seed each) are in Table~\ref{S-tab:wsfull} of the supplementary material, and the policies trained from scratch in Table~\ref{tab:viol}.}\label{fig:tradeoff}\end{figure}

\subsection{Reward penalty and constrained reinforcement learning}\label{sec:penlag}
Reward penalties and constrained reinforcement learning are the usual alternatives to a filter. The warm-started penalty policies with weights $w$ of \res{pen_ws_w_full} violate both limits at every state of charge, and their violation times decrease as $w$ increases (Table~\ref{tab:ws}). No penalty policy hopped without violating a limit (Tables~\ref{tab:viol}, \ref{tab:ws} and~\ref{S-tab:perseedws}).

The SAC-Lagrangian policies violate both limits as well, and seed 1 fell in every episode (Tables~\ref{tab:ws} and~\ref{S-tab:perseedws}). Training did not meet the constraint in any seed (Table~\ref{tab:lag}). Figure~\ref{fig:training} shows the training curves of the warm-started baseline, filter, penalty ($w=\res{figtrain_penw}$) and SAC-Lagrangian runs; the training episodes of SAC-Lagrangian seed 1 shorten during training.\begin{figure}[tbp]\centering\genfig{F7_training}{\linewidth}\caption{Training of the warm-started runs, each continued from the baseline policy of the same seed: episode return (top) and episode length in control steps (bottom), each the mean of the last \res{figtrain_window} training episodes (of all episodes so far at the start), against environment steps. Each return is in the run's own training reward: the baseline reward, the filter-mode reward with the correction penalty, the penalty-mode reward with $w=\res{figtrain_penw}$, and for SAC-Lagrangian the baseline reward without the multiplier term. Returns are therefore comparable only between the baseline and SAC-Lagrangian columns. Shaded: the first \res{sac_learnstarts} steps, in which the replay buffer is filled by the loaded policy and no gradient steps are taken.}\label{fig:training}\end{figure}
\begin{table}[tbp]\centering\caption{SAC-Lagrangian training, per seed: number of training episodes, multiplier $\mu$ at the end of training and the largest value it reached, and the mean episode cost $\bar\ell$ and episode length (Length) over the last \res{lag_last} training episodes. Cost limit $d=\res{lag_d}$; an episode has at most \res{env_maxsteps} control steps.}\label{tab:lag}
\gentable{tab_lagrangian}
\end{table}


\section{Ablation study}\label{sec:ablation}
This section varies one factor at a time: the initialisation of training, where the filter acts, the barrier rate, the margins, the controller that enforces the budget, the source of the battery parameters, and the plant parameters. Table~\ref{tab:ablsum} summarises selected settings with the same measures, and Table~\ref{S-tab:ablsumfull} of the supplementary material lists every setting; the subsections give the details. The settings other than the main evaluation use the first \res{eval_eps_set} reset seeds of each seed and state of charge.
\begin{table}[tbp]\centering\caption{Ablation summary, selected settings; policies trained with the filter unless the setting names another policy. Initialisation and filter placement: main evaluation (\res{eval_eps_main} episodes per seed and state of charge, warm-started policies for the placement); other factors: warm-started policies, \res{eval_eps_set} episodes per seed and state of charge. Violation: largest over the three states of charge of the mean over \res{n_seeds} seeds of the time below $V_{\min}$ and of the time above $I_{\lim}$; falls: count over seeds and states of charge; max height: mean over seeds at state of charge 0.2; filter active: fraction of substeps with the filter (or per-joint limit) active or infeasible, mean over states of charge. Plant variants: largest violation time and total falls over the \res{abs_nvar} variants, and ranges over them. Every setting, with the two violation times separately: Table~\ref{S-tab:ablsumfull} of the supplementary material.}\label{tab:ablsum}
\gentable{tab_ablation_summary_short}
\end{table}

\subsection{Warm start and training from scratch}\label{sec:wsscratch}
We first ask whether the results depend on the warm start. Trained from scratch, the filter-trained policies also have no violation (Table~\ref{tab:viol}), but only one of the \res{n_seeds} seeds hops repeatedly (Table~\ref{S-tab:perseed}). The comparison is not between equal training budgets: the warm-started policies add \res{train_steps_ws} steps to the \res{train_steps_scratch} steps of their baseline policy.

\subsection{Filter at test time and filter during training}\label{sec:testtime}
Next we ask whether the filter is needed during training or only at deployment. We add the filter at test time to two kinds of policy trained without it: the baseline policies, and the warm-started penalty policies with $w=\res{pst_w}$, the weight trained with \res{n_seeds} seeds that has the largest mean max height with the test-time filter at every state of charge (Table~\ref{S-tab:wsfull}). The test-time filter removes all violation and collapse time in every evaluation of both (Tables~\ref{tab:viol}, \ref{tab:ws} and~\ref{tab:abl}), but some episodes then end in a fall: the baseline policies fell in the main evaluation, both trained from scratch and warm-started (Tables~\ref{S-tab:perseed} and~\ref{S-tab:perseedws}).
\begin{table}[tbp]\centering\caption{Stress tests of the warm-started policies, one row per setting: parameter variants (plant and filter alike, except current lag and gear efficiency, plant only), barrier rate $\gamma$ (base \res{env_gamma}), and model error of the true-state filter (estimation noise $\sigma$ on its $E$ and $R_0$, or the plant's $R_0$ scaled while the filter keeps the model's; filter-trained policies only, ``--''). Penalty $w=\res{pst_w}$ + test filter: the warm-started penalty policies with weight $w$, the filter added at test time. \res{n_seeds} seeds $\times$ \res{eval_eps_set} episodes per state of charge. Violation times: largest over the states of charge of the mean over seeds (with the filter: largest of the times below $V_{\min}$, above $I_{\lim}$ and in collapse); falls: count over seeds and states of charge; filter active: mean over states of charge. Per policy, with hops, max height at every state of charge and infeasible time: Table~\ref{S-tab:ablfull} of the supplementary material.}\label{tab:abl}
\gentable{tab_ablation_short}
\end{table}

In the main evaluation neither the filter-trained policies nor the penalty policies with the test-time filter fell (Table~\ref{tab:ws}). The stress tests separate the three (Table~\ref{tab:abl}). In every parameter variant the baseline policies with the test-time filter fell in \res{abl_bft_falls_min} to \res{abl_bft_falls_max} of \res{abl_bft_neps} episodes. The penalty policies with the test-time filter fell in \res{pst_falls} of \res{pst_neps} episodes over the variants and barrier rates: seed \res{pst_fall_seed} fell in \res{pst_fall_k} of its \res{eval_eps_set} episodes at state of charge \res{pst_fall_soc} with the cold battery, with the phase resistance \res{var_rphase_hi} times its base value and with gear efficiency \res{var_geareff}, and no other episode fell. The filter-trained policies fell in none of their \res{abl_fws_neps} episodes.

Where they did not fall, the penalty policies with the test-time filter hopped higher. Their mean max height was above that of the filter-trained policies in \res{pst_nhigher} of the \res{pst_nrows} combinations of setting and state of charge, by up to \res{pst_dhmax}~m, and below it in the other \res{pst_nlower}, which are those with falls (Table~\ref{S-tab:ablfull} of the supplementary material).

\subsection{Barrier rate}\label{sec:gamma}
The barrier rate $\gamma$ sets how fast the bus may approach a limit. Changing the barrier rate from $\gamma=\res{env_gamma}$ to \res{set_gammas} leaves every filter-mode row free of violations (Table~\ref{tab:abl}). For the filter-trained policies the barrier rate barely changes the max height, and the filter acts more often at smaller rates (Table~\ref{tab:ablsum}).

\subsection{Joint-level limit and margins}\label{sec:pjm}
To separate the coupling of the drives through one budget from the budget itself, we also evaluate a joint-level limit in the spirit of the per-joint voltage-constrained bounds of Zhang et al.\ \cite{zhang2026vra}, in which each joint is limited on its own. It splits the filter's budget $B$, computed with the same barrier rate and margins, equally between the two drives. Each joint receives the current in its interval $[l_i,u_i]$ nearest to $i^{\mathrm{des}}_{q,i}$ whose power $a_i i_{q,i}+b\,i_{q,i}^2$ does not exceed $B/2$; a joint for which no such current exists receives the current of least power in its interval, as in the infeasible case of the filter. Because the two shares sum to $B$, a step in which both joints meet their shares also meets the joint budget, so the per-joint limit enforces the same bound as the filter and differs in the currents it chooses. It is applied at test time only, to the warm-started baseline and filter-trained policies. We vary the margins for the warm-started filter-trained policies: $\delta_I$ of \res{pm_dilist}~A with $\delta_V=\res{pm_mv0}$~V, $\delta_V$ of \res{pm_dvlist}~V with $\delta_I=\res{pm_mi0}$~A, and both margins zero. Table~\ref{tab:pjm} and Figure~\ref{fig:settings} give these evaluations with the base rows of the same episodes; none of its \res{pm_nrows} rows has a violation or a bus collapse.
\begin{figure}[tbp]\centering\genfig{F8_settings}{\linewidth}\caption{Joint filter, per-joint limit and margin settings, warm-started policies, \res{eval_eps_set} episodes per seed and state of charge (Table~\ref{tab:pjm}): max height against the fraction of substeps with the filter or the per-joint limit active. None of these settings has violation or collapse time. Markers: mean over seeds; bars: standard deviation over seeds; rings labelled ``fell $k/n$'': $k$ of the $n$ episodes over seeds ended in a fall.}\label{fig:settings}\end{figure}

\emph{Per-joint limit.} The baseline policies fell in more episodes with the per-joint limit at test than with the joint filter at test; the filter-trained policies did not fall with either, but hopped lower with the per-joint limit (Table~\ref{tab:pjm} and Figure~\ref{fig:settings}).

\emph{Margins.} The current margin sets the mean per-episode peak battery current to $I_{\lim}-\delta_I$ (Table~\ref{tab:pjm}); with $\delta_I=0$ the battery current reaches $I_{\lim}=\res{pm_ilim}$~A, but no substep exceeds it. The voltage margin has almost no effect on max height and filter activity (Table~\ref{tab:pjm}), consistent with the current limit binding first (Section~\ref{sec:stress}). With both margins zero, in \res{pm_deltaV0.0deltaI0.0_filtws_neps} episodes no limit was breached, the bus did not collapse, no episode fell and no substep was infeasible. This is consistent with Proposition~\ref{prop:budget} and Corollary~\ref{cor:closed}, which need no margin when the battery parameters are exact; the margins are there for errors in the battery model (Remark~\ref{rem:sens}).
\begin{table}[tbp]\centering\caption{Joint-level limit and margins, warm-started policies, \res{n_seeds} seeds $\times$ \res{eval_eps_set} episodes per state of charge (the episodes of the stress tests). Per-joint limit: the filter's budget split equally between the two drives, each joint limited on its own, at test time only. Margins: the policies trained with the filter, evaluated with the joint filter and the margins given. Violation and collapse times: largest over the three states of charge of the mean over seeds; peak $\Ib$: largest over the states of charge of the mean per-episode peak; infeasible time likewise. Falls: count over seeds and states of charge. Hops and filter or limiter active: mean over states of charge.}\label{tab:pjm}
\gentable{tab_perjoint_margin}
\end{table}

\subsection{Estimated battery parameters}\label{sec:estresults}
The guarantee assumes known $E$ and $R_0$; this test drops that assumption. Table~\ref{tab:est} compares the filter-trained policies under the true-state filter and the EKF filter, with fixed or confidence-scaled margins (Section~\ref{sec:estimator}). No episode ends in a fall (\res{est_falls_total} of \res{est_neps_total}), and the proof-assumption counters sum to \res{est_counters_total} substeps over all rows.

\emph{Nominal conditions.} Without model or sensor error, the EKF filter matches the true-state filter.

\emph{Plant resistance \res{set_rsmall} times the model's.} The true-state filter, which keeps the model's $R_0$, breaches the current limit at every state of charge; the EKF filter, with either margin, removes these breaches. The EKF learns the plant's resistance: in all conditions without injected noise its final estimate of $R_0$ is \res{est_clean_r0rel_min} to \res{est_clean_r0rel_max} times the plant's on average.

\emph{Plant resistance \res{set_rbig} times the model's.} Here the true-state filter stays within the limits but is conservative: the EKF filter draws a larger peak battery current and hops higher.

\emph{Mis-specified sensor noise ($\sigma=\res{est_sigma}$), a stress test.} Here the EKF fails. Each EKF measurement is multiplied by $1+\sigma n$, with an independent standard normal $n$ at every substep. The true-state filter receives a different error, a perturbation of $E$ and $R_0$. The injected voltage noise has a standard deviation of \res{est_noisev_lo} to \res{est_noisev_hi}~V at the open-circuit voltage of the evaluated states of charge, \res{est_ratiov_lo} to \res{est_ratiov_hi} times the EKF's assumed $\sigma_v=\res{par_sigv}$~V. The injected current noise at \res{est_itarget}~A is \res{est_noisei}~A, \res{est_ratioi} times its assumed $\sigma_i=\res{par_sigi}$~A. Under this mis-specification the EKF's confidence is \res{est_noise_confmax} in every episode, its final $R_0$ estimate diverges in \res{est_noise_rdiv} of \res{est_noise_neps} episodes (Table~\ref{tab:est}), and its mean $|\hat E - E|$ is \res{est_noise_eabs_min} to \res{est_noise_eabs_max}~V, against at most \res{est_clean_eabs_max}~V without injected noise. With fixed margins both filters breach the current limit (Table~\ref{tab:est}). Confidence-scaled margins reduce the breaches but do not remove them; they remain at \res{est_noise_ekfconf_nsococ} state of charge (1.0).

\emph{Confidence acted as a switch.} The episode-mean confidence was at least \res{est_clean_confmin} in every episode without injected noise and \res{est_noise_confmax} in every episode with it, so the margins \eqref{eq:confmargin} were at their fixed values or at \res{env_marginratio} times those values. These runs cannot distinguish confidence-scaled margins from a fixed \res{env_marginratio}$\times$ margin switched on when the estimator fails.
\begin{table}[tbp]\centering\caption{Estimated battery parameters: warm-started filter policies with the true-state filter, the EKF with fixed margins and the EKF with confidence-scaled margins, under nominal conditions, the plant's $R_0$ scaled, and mis-specified measurement noise $\sigma$ (true state: on $E$ and $R_0$; EKF: on its two measurements). \res{n_seeds} seeds $\times$ \res{eval_eps_set} episodes per state of charge. $t(\Ib>I_{\lim})$: largest over the states of charge of the mean over seeds (no row has time below $V_{\min}$); largest $\Ib$: over all episodes; falls: over seeds and states of charge; max height: mean over seeds; confidence: smallest to largest episode mean; $R_0$ diverged: episodes whose final $R_0$ estimate is outside \res{est_rdiv_lo} to \res{est_rdiv_hi} times the plant's. Per state of charge, with standard deviations and the proof-assumption counters: Table~\ref{S-tab:estfull} of the supplementary material.}\label{tab:est}
\gentable{tab_estimator_short}
\end{table}

\subsection{Plant variants and model error}\label{sec:stress}
Finally we test the filter on plants that differ from its model. Table~\ref{tab:abl} gives the warm-started policies under the parameter variants, the barrier rates and the model errors. For each of the \res{abl_fm_ngrp} groups of policies run with the filter (the filter-trained policies, and the baseline and penalty policies with the test-time filter), the violation and collapse times are zero in all \res{abl_fm_nset} settings of the variants and barrier rates. This includes the current-loop lag and the gear efficiency, which only the plant has. The baseline policies without the filter violate a limit in all \res{abl_nvar} variants.

Lowering $V_{\min}$ from \res{par_vmin}~V to the cell datasheet's \res{par_vmin_ds}~V leaves the filter rows essentially unchanged, because the current limit binds first (Table~\ref{tab:abl}, and Table~\ref{S-tab:vminid} of the supplementary material).

With a wrong battery model, the true-state filter breaches the current limit in \res{abl_me_nviol} of the \res{abl_me_n} model-error settings (Table~\ref{tab:abl}). Estimation noise ($\sigma=\res{set_noise_lo}$ and \res{set_noise_hi}) gives the breaches analysed in Remark~\ref{rem:sens} and Table~\ref{tab:sens}. When the plant's $R_0$ is \res{set_rsmall} times the model's, the battery current exceeds $I_{\lim}$ at every state of charge. This is the direction that Remark~\ref{rem:sens} predicts to be unsafe: with $\hat R_0 > R_0$ the filter predicts a smaller current than the plant draws at the same end-of-step voltage. When the plant's $R_0$ is larger (\res{set_rmid} and \res{set_rbig} times), no limit is violated, but the policies hop lower.

